\chapter{RC Circuits}
\label{ch:rc}
RC circuits appear as coupling and bypass networks, timing circuits, audio
filters, and detector components. An RC circuit is the electrical embodiment of
the \emph{first-order system} of \cref{ch:linsys}: one energy store, one time
constant, one real pole.

\section{The Circuit}
Take a source \(u(t)\) driving a resistor \(R\) in series with a capacitor \(C\)
(\cref{fig:rc-schematic}). The natural state variable is the capacitor voltage
\(v_C\), because it holds the stored energy.

We must also say what the \emph{output} is, since one circuit can deliver several.
Write \(v_o\) for the output voltage---the quantity we measure and the one the
transfer function refers to. Here we take it across the capacitor, so
\[
v_o(t)=v_C(t),
\]
which turns out to be the low-pass response. Taking \(v_o\) across the resistor
instead, \(v_o=v_R=u-v_C\), gives the high-pass response from the very same
circuit; \cref{sec:rc-freq} works out both. Whenever a circuit's behavior is
described in this book, it is behavior \emph{from \(u\) to \(v_o\)}, and the choice
of \(v_o\) is part of the model.

\begin{figure}[htbp]
\centering
\begin{circuitikz}[scale=1]
\draw (0,0) to[sV,l=\(u(t)\)] (0,2)
      to[R,l=\(R\),i>=\(i(t)\)] (3,2)
      to[C,l=\(C\),v=\(v_C(t)\)] (3,0) -- (0,0);
\draw (3,2) to[short,-o] (4,2) node[right]{$+$};
\draw (3,0) to[short,-o] (4,0) node[right]{$-$};
\node at (4.4,1) {\(v_o\)};
\end{circuitikz}
\caption{Series RC circuit. The input is the source \(u(t)\); the output \(v_o\)
is taken across the capacitor here, so \(v_o=v_C\) and the circuit is a low-pass
filter. Taking \(v_o\) across the resistor instead gives a high-pass filter.}
\label{fig:rc-schematic}
\end{figure}

\section{From Schematic to Differential Equation}
KVL around the loop says the source equals the resistor drop plus the capacitor
voltage, and the capacitor law relates its current to its rate of voltage change:
\[
u(t)=R\,i(t)+v_C(t),
\qquad
i(t)=C\frac{\dd v_C}{\dd t}.
\]
Substituting the second into the first,
\[
\boxed{\,RC\frac{\dd v_C}{\dd t}+v_C=u(t).\,}
\]

\section{State-Space Form and the Transfer Function}
With the single state \(x=v_C\), and output \(y=v_o=v_C\),
\[
\dot{x}=-\frac{1}{RC}\,x+\frac{1}{RC}\,u,
\qquad
y=x=v_o .
\]
There is one state because there is one independent energy store, so there is one
pole. To make that pole explicit, Laplace-transform either the boxed ODE or the
state equation---both give the same thing. From the state equation with zero
initial conditions,
\[
sX=-\frac{1}{RC}X+\frac{1}{RC}U
\;\Rightarrow\;
\Bigl(s+\frac{1}{RC}\Bigr)X=\frac{1}{RC}U
\;\Rightarrow\;
G(s)=\frac{Y}{U}=\frac{X}{U}=\frac{1/RC}{s+1/RC}=\frac{1}{1+RCs}.
\]
Now the pole is visible as the value of \(s\) that makes the denominator vanish:
\[
s_{\text{pole}}=-\frac{1}{RC}.
\]
The two forms say the same thing---the state-space coefficient \(-1/RC\) \emph{is}
the pole, because for a single state the \(1\times1\) ``matrix'' is its own
eigenvalue---but the transfer function displays it directly.

\section{Matching the Generic First-Order System}
Compare with the canonical first-order system \(\tau\dot y+y=Ku\) from
\cref{ch:linsys}. They are identical with
\[
\tau=RC,
\qquad
K=1,
\qquad
s_{\text{pole}}=-\frac{1}{\tau}=-\frac{1}{RC}.
\]
Every first-order result from \cref{ch:linsys} now applies with no extra work:
the step response, the \(63.2\%\) rule, the corner frequency, and the Bode shape.

\section{Transient Response}
For a DC step of height \(E\) applied at \(t=0\) with the capacitor initially
uncharged, the two outputs are
\[
v_C(t)=E\left(1-\ee^{-t/RC}\right),
\qquad
v_R(t)=E-v_C(t)=E\,\ee^{-t/RC},
\]
and a capacitor discharging from an initial voltage \(V_0\) with the source removed
follows \(v_C(t)=V_0\,\ee^{-t/RC}\). These are the first-order responses of
\cref{ch:linsys} read as circuit voltages (\cref{fig:rc-transient}): after one time
constant \(\tau=RC\) the capacitor has charged to \(63.2\%\) of \(E\) while the
resistor voltage has fallen to \(36.8\%\), and after about \(5\tau\) the circuit is
settled.

\begin{figure}[htbp]
\centering
\includegraphics[width=0.82\textwidth]{rc_transient}
\caption{RC step response, both outputs. The capacitor voltage \(v_C\) rises
toward \(E\) while the resistor voltage \(v_R=E-v_C\) decays---mirror exponentials
sharing the time constant \(\tau=RC\), one \(\tau\) covering \(63.2\%\) of the
change. These are the time-domain faces of the low-pass and high-pass responses of
\cref{sec:rc-freq}.}
\label{fig:rc-transient}
\end{figure}

\section{Frequency Response: Low-Pass and High-Pass}
\label{sec:rc-freq}
The same circuit is a low-pass or a high-pass filter depending on where \(v_o\) is
taken, and both transfer functions come from the one state-space model already in
hand---no new circuit analysis is needed.

\emph{Output across \(C\).} This is the case developed above, \(y=v_C\), so
\(G_C(s)=1/(1+RCs)\) directly.

\emph{Output across \(R\).} The output equation changes but the state equation does
not. KVL gives \(v_R=u-v_C=u-x\), so in state-space terms only \(C_{\text{out}}\)
and \(D\) change:
\[
\dot x=-\frac{1}{RC}x+\frac{1}{RC}u,
\qquad
y=v_R=\underbrace{(-1)}_{C_{\text{out}}}x+\underbrace{(1)}_{D}u .
\]
Transforming, \(Y=-X+U\) with \(X=G_C(s)\,U\) from before:
\[
G_R(s)=\frac{Y}{U}=1-G_C(s)=1-\frac{1}{1+RCs}=\frac{RCs}{1+RCs}.
\]
So the two responses are complements, \(G_C+G_R=1\)---at every frequency, whatever
the capacitor does not pass, the resistor does. On the imaginary axis,
\[
G_C(\jj\omega)=\frac{1}{1+\jj\omega RC}\ \text{(low-pass)},
\qquad
G_R(\jj\omega)=\frac{\jj\omega RC}{1+\jj\omega RC}\ \text{(high-pass)}.
\]
The \(D\ne0\) term in the second model is worth noticing: it is why the high-pass
has a \emph{zero} (at \(s=0\), from the factor \(RCs\)) while the low-pass has none.
A direct input-to-output path creates a zero.

\begin{workedbox}[: which output, and at what frequency?]
Take \(R=\SI{1}{\kilo\ohm}\) and \(C=\SI{160}{\pico\farad}\), so
\(\tau=\SI{160}{\nano\second}\) and \(f_c=1/(2\pi\tau)\approx\SI{995}{\kilo\hertz}\)
---call it \SI{1}{\mega\hertz}. At \SI{100}{\kilo\hertz}, one decade \emph{below}
the corner, \(\omega\tau=0.1\) and
\[
|G_C|=\frac{1}{\sqrt{1+0.01}}\approx0.995\ \ (-\SI{0.04}{\decibel}),
\qquad
|G_R|=\frac{0.1}{\sqrt{1+0.01}}\approx0.0995\ \ (-\SI{20}{\decibel}).
\]
At \SI{10}{\mega\hertz}, one decade \emph{above}, the two swap exactly. And at the
corner itself both are \(1/\sqrt2\)---the one frequency where the circuit splits
the signal evenly, which is what \(|G_C|^2+|G_R|^2=1\) demands. Same components,
same instant: the only thing that changed is which element you measured across.
\end{workedbox}

Both share the corner frequency \(\omega_c=1/(RC)\), where the gain is
\(-\SI{3}{\decibel}\). The low-pass passes slow signals and rolls off high ones
(\cref{fig:bode1}); the high-pass does the reverse (\cref{fig:rc-highpass}).

\begin{figure}[htbp]
\centering
\includegraphics[width=0.82\textwidth]{bode_first_order_highpass}
\caption{RC high-pass response (output across \(R\)): \(-\SI{20}{\decibel}\) per
decade rise below the corner, flat above, phase from \(+90^\circ\) to
\(0^\circ\).}
\label{fig:rc-highpass}
\end{figure}

\section{The Pole and the Root Locus}
The single pole sits on the negative real axis at \(-1/(RC)\)
(\cref{fig:rc-rootlocus}). Changing \(R\) or \(C\) slides it along that axis: a
larger \(RC\) (slower circuit, lower cutoff) moves the pole toward the origin; a
smaller \(RC\) moves it left toward faster response. Because the pole never leaves the real axis,
an RC circuit cannot overshoot or ring---the first-order signature.

\begin{figure}[htbp]
\centering
\includegraphics[width=0.78\textwidth]{rootlocus_first_order}
\caption{Root locus of the RC pole. The single real pole \(-1/(RC)\) slides along
the negative real axis as the time constant changes; it never becomes complex, so
there is no ringing.}
\label{fig:rc-rootlocus}
\end{figure}

\begin{controlsbox}
\textbf{Connection.} The RC circuit is the textbook first-order lag. Its cutoff
frequency, its \(-\SI{20}{\decibel}\)/decade roll-off, and its \(63.2\%\)-per-time-constant
step are three faces of the single pole at \(-1/\tau\). Cascading RC stages adds
poles and steepens the roll-off, at the cost of accumulated phase lag---the same
bookkeeping as adding first-order lags in a control loop.
\end{controlsbox}

\section{Energy Interpretation}
The capacitor stores \(E_C=\tfrac12 Cv_C^2\). During charging, source energy
splits between the electric field and resistive heat; during discharge, the
stored field energy is dissipated in \(R\).

\begin{energybox}
The time constant \(\tau=RC\) is how long that energy exchange takes. Over a full
charge--discharge cycle the capacitor returns everything it stored, so the resistor is
the only place energy actually leaves---the time-domain reason a first-order circuit
has nothing to sustain an oscillation with.
\end{energybox}

\begin{unitsbox}
\[
RC=\si{\second},\qquad
\omega_c=\frac{1}{RC}=\si{\per\second},\qquad
f_c=\frac{1}{2\pi RC}=\si{\hertz}.
\]
\end{unitsbox}

\begin{mistakebox}
A capacitor is an open circuit only at DC steady state. During a transient or at
AC, its current \(i=C\,\dd v_C/\dd t\) is generally nonzero.
\end{mistakebox}

\begin{exambox}
Three numbers answer almost every RC question, and all three are the single pole
\(-1/RC\) wearing different clothes: the time constant \(\tau=RC\) (one \(\tau\)
covers \SI{63.2}{\percent} of a step, five \(\tau\) settles it), the corner
frequency \(f_c=1/(2\pi RC)\), and the \(-\SI{20}{\decibel}\)/decade roll-off. Note
the \(2\pi\): questions give frequency in hertz and time constants in seconds, and
mixing \(f\) with \(\omega\) is the most common way to be wrong by a factor of
\(6.28\). Then ask \emph{which element the output is across}---capacitor means
low-pass, resistor means high-pass, and nothing else about the circuit changes.
\end{exambox}

\section{Worked Example}
Let \(R=\SI{4.7}{\kilo\ohm}\) and \(C=\SI{100}{\nano\farad}\). Then
\[
\tau=RC=(4.7\times10^3)(100\times10^{-9})=\SI{4.70e-4}{\second},
\qquad
f_c=\frac{1}{2\pi\tau}\approx\SI{339}{\hertz}.
\]
The pole is at \(-1/\tau\approx-\SI{2130}{\per\second}\).

\Cref{fig:rc-worked} plots this specific circuit, with real frequency in hertz on
the axis rather than the normalized axes of the generic figures. (Those generic
plots---\cref{fig:bode1,fig:rc-rootlocus}---show the \emph{shape}, which is the
same for every first-order system; this one shows \emph{this} circuit.) Reading it:
signals well below \SI{339}{\hertz} pass nearly unattenuated, signals a decade above
(\(\sim\!\SI{3.4}{\kilo\hertz}\)) are down about \SI{20}{\decibel}, and the phase
lag is \(45^\circ\) right at the corner. In the time domain the same circuit settles
a step in roughly \(5\tau\approx\SI{2.4}{\milli\second}\).

\begin{figure}[htbp]
\centering
\includegraphics[width=\textwidth]{rc_worked}
\caption{The worked example, \(R=\SI{4.7}{\kilo\ohm}\) and
\(C=\SI{100}{\nano\farad}\). \emph{Left:} its Bode plot on a real frequency axis;
the corner sits at \(f_c=\SI{339}{\hertz}\), where the gain is \(-\SI{3}{\decibel}\)
and the phase \(-45^\circ\). \emph{Right:} the root locus as \(R\) is swept with
\(C\) fixed. The pole \(s=-1/RC\) only slides along the real axis---toward the
origin (slower, lower cutoff) as \(R\) grows---so no value of \(R\) can make this
circuit ring.}
\label{fig:rc-worked}
\end{figure}

\begin{physicalbox}
\textbf{How does changing \(R\) change the dynamics?} Only one thing moves: the
pole, and only along the real axis. Doubling \(R\) doubles \(\tau\), halves \(f_c\),
and halves the pole's distance from the origin---the circuit gets uniformly slower
without any change in the \emph{character} of the response. There is no overshoot to
gain and no ringing to provoke, because a single real pole has no imaginary part to
acquire. That is the practical meaning of ``first order,'' and it is why RC timing
and coupling networks are so predictable.
\end{physicalbox}
